\section{Two surface tests for nefness}\label{num:section}

Our goal is a numerical criterion for nefness. The first surface will
convert positive intersection square into section growth with one fixed
twist. The second surface will show that a negative curve forces that
square to be positive. Together these statements let the numerical-dimension
hypothesis, imposed later on the original variety, rule out every negative
curve. The arguments apply to rational divisors on an
\(n\)-dimensional variety, for every \(n\geq 3\).  We use the following
precise approximation hypothesis:
\begin{equation}\label{num:approximation}
\begin{gathered}
Y\text{ is a normal projective }n\text{-fold over }\mathbb C,
\quad n\geq3,
\qquad \operatorname{codim}_Y\operatorname{Sing}Y\geq 3,\\
L,B_Y\text{ are }\mathbb Q\text{-Cartier},\qquad
t_j\in\mathbb Q_{>0},\quad t_j\longrightarrow 0,
\qquad M_j=L+t_jB_Y,\\
\text{for each }j\text{ there are a closed }Z_j\subset Y,
\quad\operatorname{codim}_Y Z_j\geq 3,
\quad\text{and }d_j\in\mathbb Z_{>0}\text{ such that}\\
d_jM_j\text{ is Cartier and }|d_jM_j|
\text{ is generated by global sections on }Y\setminus Z_j.
\end{gathered}
\end{equation}
Every positive multiple of \(d_j\) has the same generation property.
The integers \(d_j\) may depend on \(j\).
Restrictions of rational Cartier divisors below mean restrictions as
rational line bundles.  Neither effectiveness nor positivity of \(B_Y\) is part
of \eqref{num:approximation}.

\subsection{A fixed twist detects positive surface square}

\begin{lemma}[Surface growth with a fixed twist]\label{num:growth}
Assume \eqref{num:approximation}, and let \(H\) be a very ample Cartier
divisor on \(Y\).  Fix a projective resolution
\(\pi:\widehat Y\to Y\) which is an isomorphism over the smooth locus.
There are a smooth complete intersection surface
\(S\subset Y\), cut out by \(n-2\) members of \(|H|\), and a Cartier
divisor \(P\) on \(\widehat Y\) with the following properties.
Identify \(S\) with its inverse image in \(\widehat Y\).  Then
\(L|_S\) is nef, so
\begin{equation}\label{num:nonnegative-square}
L^2\cdot H^{n-2}\geq 0.
\end{equation}
If \(dL\) is Cartier, then for every positive integer \(m\) divisible
by \(d\) the restriction map
\begin{equation}\label{num:restriction-surjective}
H^0(\widehat Y,m\pi^*L+P)
\longrightarrow H^0(S,mL|_S+P|_S)
\end{equation}
is surjective.  In particular, as \(m\to\infty\) through multiples of
this one fixed integer \(d\),
\begin{equation}\label{num:quadratic-growth}
h^0(\widehat Y,m\pi^*L+P)
\geq \frac{m^2}{2}\,L^2\cdot H^{n-2}+O(m),
\end{equation}
where the error term depends only on the fixed surface and its divisors.
\end{lemma}

\begin{proof}
Put \(r=n-2\).  Very general members of \(|H|\) have a smooth
complete intersection \(S\) avoiding \(\operatorname{Sing}Y\) and
every \(Z_j\).  Indeed, each of these countably many closed sets has
dimension at most \(n-3\), so the condition that \(r\) hyperplanes
meet any one of them is a proper closed incidence condition.
Over \(\mathbb C\) these can be avoided simultaneously, together
with the closed conditions excluded by Bertini's theorem.
The inverse image of \(S\) is therefore exactly the smooth
complete intersection cut out by the \(r\) pulled-back equations;
there are no additional components over the singular locus.
Each \(M_j|_S\) is semiample, hence nef.  Passing to the limit on
every curve on \(S\) shows that \(L|_S\) is nef and proves
\eqref{num:nonnegative-square}.

Write \(A=\pi^*H\), choose an ample Cartier divisor \(A_0\) on
\(\widehat Y\), and make the fixed choice
\begin{equation}\label{num:fixed-twist}
P=K_{\widehat Y}+rA+A_0.
\end{equation}
Since \(A\) is nef and ampleness is open
\cite[Corollary~1.4.10 and Theorem~1.4.23]{Laz04I}, there exists \(\eta>0\)
such that
\begin{equation}\label{num:ample-errors}
A_0+(r-i)A-\epsilon\pi^*B_Y
\quad\text{is ample whenever }0\leq i\leq r
\text{ and }0\leq\epsilon<\eta.
\end{equation}
For a given positive \(m\) with \(d\mid m\), first choose
\(j=j(m)\) so that \(mt_j<\eta\), and then choose a multiple
\(k\) of \(d_j\) with \(k>m\).  The pullback of
\(|kM_j|\) is generated near \(S\).  A general member \(D_m\)
is smooth there, by Bertini, and
\[
R_m=\frac{m}{k}D_m\geq 0,
\qquad R_m\sim_{\mathbb Q}m\pi^*M_j,
\qquad \mathcal J(R_m)=\mathcal O_{\widehat Y}
\quad\text{near }S.
\]
Here \(\mathcal J(R_m)\) is the multiplier ideal; its asserted
local triviality follows from \(m/k<1\) and the smoothness of
\(D_m\) there.

The divisors \(m\pi^*L+P-iA\) are integral Cartier divisors, and
\[
m\pi^*L+P-iA-K_{\widehat Y}-R_m
\sim_{\mathbb Q}
A_0+(r-i)A-mt_j\pi^*B_Y
\]
is ample for \(0\leq i\leq r\).  Nadel vanishing
\cite[Theorem~9.4.8]{Laz04II}, equivalently
\cite[Theorem~2.4]{LazMI}, gives
\begin{equation}\label{num:nadel}
H^q\!\left(\widehat Y,
\mathcal O_{\widehat Y}(m\pi^*L+P-iA)
\otimes\mathcal J(R_m)\right)=0
\qquad(q>0,\ 0\leq i\leq r).
\end{equation}

Set \(\mathcal F=
\mathcal O_{\widehat Y}(m\pi^*L+P)\otimes\mathcal J(R_m)\).
Tensoring the Koszul complex of the \(r\) equations with
\(\mathcal F\) yields the exact complex
\[
0\longrightarrow\mathcal F(-rA)
\longrightarrow\cdots\longrightarrow
\mathcal F(-A)^{\oplus r}\longrightarrow\mathcal F
\longrightarrow\mathcal O_S(mL|_S+P|_S)
\longrightarrow 0.
\]
At points of \(S\), the equations form a regular
sequence and \(\mathcal F\) is locally free.  Outside \(S\), one
equation is a unit, so the Koszul complex is contractible even after
tensoring with an arbitrary sheaf.  This proves exactness everywhere.
Splitting the complex into short exact sequences and using
\eqref{num:nadel} gives surjectivity on global sections onto its last
term.  The inclusion
\(\mathcal F\subset\mathcal O_{\widehat Y}(m\pi^*L+P)\)
then proves \eqref{num:restriction-surjective}.

Adjunction and \eqref{num:fixed-twist} give
\(P|_S-K_S=A_0|_S\).  Hence \(mL|_S+P|_S-K_S\) is ample.
Kodaira vanishing, the zero-boundary ample case of
\cite[Theorem~2.4]{LazMI}, and surface Riemann--Roch
\cite[Chapter~V, Section~1]{Hartshorne77} give
\[
\begin{split}
h^0(S,mL|_S+P|_S)
={}&\frac{m^2}{2}(L|_S)^2
+\frac m2 L|_S\cdot(2P|_S-K_S)\\
&+\frac12 P|_S\cdot(P|_S-K_S)
+\chi(S,\mathcal O_S).
\end{split}
\]
If \(S\) is disconnected, the formula is read componentwise and
summed.  All its coefficients are fixed before \(m,j,k,D_m\) vary.
Together with \eqref{num:restriction-surjective}, this proves
\eqref{num:quadratic-growth}.
\end{proof}

\subsection{Negative degree forces ambient vanishing}

Lemma~\ref{num:growth} supplies the first surface test. To obtain
nefness once its square is zero, we must rule out curves missed by that
very general surface, including curves in the singular locus. The next
two lemmas will measure their effect on a second surface.

The next elementary estimate retains the cotangent bundle of the
ambient variety.  This is what permits a curve singular at the
point where vanishing will be measured.

\begin{lemma}[Uniform ambient multiplicity]\label{num:jets}
Let \(V\) be a smooth projective complex variety, let
\(\Gamma\subset V\) be an integral curve, and let
\(\iota:N\to V\) be the map from its normalization.
For a section \(s\), write \(\operatorname{ord}_z(s)\) for its order
in the regular local ring of \(V\) at \(z\).
There is a positive constant \(g\), depending only on this map,
such that for every line bundle \(\mathcal D\) on \(V\), every
nonzero section \(s\in H^0(V,\mathcal D)\), and every closed
point \(z\in\Gamma\),
\begin{equation}\label{num:jet-bound}
\operatorname{ord}_z(s)\geq
-\frac{\deg_N\iota^*\mathcal D}{g}.
\end{equation}
\end{lemma}

\begin{proof}
Choose a line bundle \(G\) of positive degree on the fixed smooth
curve \(N\) such that
\(\iota^*\Omega^1_V\) embeds as a subbundle of
\(G^{\oplus a}\) for some \(a\).  Such a choice follows by
global generation of
\((\iota^*\Omega^1_V)^*\otimes G\) for a sufficiently ample
\(G\), followed by dualization.  Set \(g=\deg G>0\).

For a nonzero \(s\), the integer
\[
b=\min_{p\in\Gamma\text{ closed}}\operatorname{ord}_p(s)
\]
is finite and attained.  The leading ambient jet gives a nonzero
section of
\begin{equation}\label{num:leading-jet}
\operatorname{Sym}^b(\iota^*\Omega^1_V)
\otimes\iota^*\mathcal D.
\end{equation}
Here is a justification valid also at singular points of \(\Gamma\).
For \(b\geq1\), the bundles of principal parts on the smooth
variety \(V\) have the locally split exact sequence
\[
0\longrightarrow\operatorname{Sym}^b\Omega^1_V\otimes\mathcal D
\longrightarrow\mathcal P^b(\mathcal D)
\longrightarrow\mathcal P^{b-1}(\mathcal D)
\longrightarrow 0.
\]
The lower jet of \(s\) has zero value at every closed point of
the reduced curve \(\Gamma\). A section of a vector bundle on a
reduced variety with this property is zero, so that lower jet restricts
to the zero section. The principal-parts sequence is locally split
as a sequence of \(\mathcal O_V\)-modules, and therefore remains
exact after restriction to \(\Gamma\) and pullback to \(N\).
The restricted order-\(b\) jet consequently belongs to the left-hand
bundle. Its value is nonzero at a point where the minimum \(b\)
is attained; the induced map on its fiber at any point of \(N\)
above that point is an isomorphism. Its pullback is thus regular and
nonzero.  For \(b=0\), use \(s|_\Gamma\)
directly.  This proves \eqref{num:leading-jet} without assuming
smoothness of \(\Gamma\).

Taking symmetric powers of the subbundle inclusion above embeds
\eqref{num:leading-jet} into a direct sum of copies of
\(G^{\otimes b}\otimes\iota^*\mathcal D\).  A nonzero component
is a nonzero section of this line bundle on \(N\), so
\[
0\leq bg+\deg_N\iota^*\mathcal D.
\]
Since \(\operatorname{ord}_z(s)\geq b\) for every \(z\in\Gamma\),
\eqref{num:jet-bound} follows.
\end{proof}

\subsection{A fixed exceptional curve forces a positive square}

\begin{lemma}[Surface fixed-part estimate]\label{num:fixed-square}
Let \(T\) be an integral projective surface, let
\(f:T_1\to T\) be a projective birational morphism from a smooth
surface, and fix a curve \(e\subset T_1\) contracted by \(f\).
For each \(j\), let \(D_j\) be a rational Cartier divisor on
\(T\), and let \(k_j>0\) be an integer such that \(k_jD_j\)
is Cartier.  Suppose a nonzero linear system of
\(k_jf^*D_j\) is generated outside the inverse image of a finite
subset of \(T\).  Let \(F_j\) be its fixed divisor.
If, for one fixed \(c>0\),
\[
\operatorname{coeff}_e F_j\geq ck_j
\quad\text{for all }j,
\]
then there exists \(\epsilon>0\), independent of \(j\), such that
\begin{equation}\label{num:surface-positive-square}
D_j^2\cdot[T]\geq\epsilon
\quad\text{for all }j.
\end{equation}
No normality assumption on \(T\) is needed.
\end{lemma}

\begin{proof}
Every curve in \(\operatorname{Supp}F_j\) maps to a point of
\(T\), by the generation assumption.  Fix a very ample Cartier
divisor \(H_T\) on \(T\) and an ample Cartier divisor \(J\) on
\(T_1\), and put \(Q=f^*H_T\).  The projection formula gives
\begin{equation}\label{num:orthogonality}
Q^2=H_T^2\cdot[T]>0,
\qquad Q\cdot F_j=f^*D_j\cdot F_j=0.
\end{equation}
The Hodge index theorem on the fixed smooth surface \(T_1\)
\cite[Chapter~V, Section~1]{Hartshorne77} makes the negative intersection
form positive definite on
\(Q^\perp\subset N^1(T_1)_{\mathbb R}\).  Consequently there
is one constant \(K>0\) such that
\[
(J\cdot F)^2\leq K(-F^2)
\qquad\text{for every }F\in Q^\perp.
\]
Effectiveness of \(F_j\) gives
\(J\cdot F_j\geq ck_j(J\cdot e)\).  Thus, with the fixed
positive number \(\epsilon=c^2(J\cdot e)^2/K\),
\begin{equation}\label{num:fixed-negative-square}
F_j^2\leq-\epsilon k_j^2.
\end{equation}

After removal of \(F_j\), the residual linear system has no
fixed curve.  Its square is nonnegative: choose two members
with no common curve and take their effective intersection,
or use a nowhere vanishing member if the system is trivial.
By \eqref{num:orthogonality} and the projection formula,
\[
0\leq(k_jf^*D_j-F_j)^2
=k_j^2(D_j^2\cdot[T])+F_j^2.
\]
Now \eqref{num:fixed-negative-square} proves
\eqref{num:surface-positive-square}.
\end{proof}

\subsection{The nefness criterion}

We now combine the last two estimates. A negative curve forces all
sections to vanish linearly at one fixed ambient point. Restriction to
a surface through that point, followed by one blowup, produces the fixed
exceptional curve required by Lemma~\ref{num:fixed-square}. The choices
are made before the perturbation index varies.

\begin{proposition}[Vanishing surface square implies nefness]
\label{num:nef}
Assume \eqref{num:approximation}.  If, for one very ample Cartier
divisor \(H\) on \(Y\),
\begin{equation}\label{num:zero-square}
L^2\cdot H^{n-2}=0,
\end{equation}
then \(L\) is nef.
\end{proposition}

\begin{proof}
Suppose that an integral curve \(C_0\subset Y\) satisfies
\(L\cdot C_0<0\).  Fix a projective resolution
\(\pi:\widehat Y\to Y\) which is an isomorphism over the smooth
locus.  Choose a closed point \(x\in C_0\) outside the images
of those irreducible components of \(\pi^{-1}(C_0)\) which do
not dominate \(C_0\).  There are only finitely many such images,
and each is a point.  Consequently every point of \(\pi^{-1}(x)\)
lies on a component of \(\pi^{-1}(C_0)\) dominating \(C_0\).

\smallskip\noindent
\emph{A surface through \(x\).}
Choose \(n-2\) very general members
\(H_1,\ldots,H_{n-2}\in|H|\), all passing through \(x\).
They meet properly, and their intersection meets each \(Z_j\)
and \(\operatorname{Sing}Y\) in at most finitely many points.
Indeed, the hyperplanes through \(x\) have no other base point;
successive general cuts avoid all positive-dimensional components
of the relevant intersections.  The countably many conditions
can again be imposed simultaneously over \(\mathbb C\).
Write their intersection cycle as
\[
H_1\cdots H_{n-2}\cdot[Y]=\sum_{i=1}^a a_i[T^{(i)}],
\qquad a_i>0,
\]
where each \(T^{(i)}\) is an integral surface.
The dimension theorem guarantees a component through \(x\);
fix one and denote it by \(T\).

On every \(T^{(i)}\), the system induced by \(|d_jM_j|\) has at
most a finite base locus.  Its square is nonnegative, including
when \(T^{(i)}\) is nonnormal.  To check this directly, choose a
section nonzero at the generic point, and then choose a second
section not vanishing identically on any curve of the first
zero divisor.  The resulting Cartier intersections with
\([T^{(i)}]\) form an effective zero-cycle.  If the first section
has no zero, the restricted line bundle is trivial and its
square is zero.  These are the usual Cartier intersection and
projection formulas on cycles \cite[Chapter~2]{Fulton98}.
It follows that
\[
M_j^2\cdot[T^{(i)}]\geq0,
\qquad
\sum_{i=1}^a a_i(M_j^2\cdot[T^{(i)}])
=M_j^2\cdot H^{n-2}\longrightarrow 0
\]
by \eqref{num:zero-square}.  In particular,
\begin{equation}\label{num:special-surface-limit}
M_j^2\cdot[T]\longrightarrow 0.
\end{equation}

\smallskip\noindent
\emph{Fixed geometry above \(x\).}
The surface \(T\) meets \(\operatorname{Sing}Y\) only in finitely
many points.  Its strict transform \(T_{\widehat Y}\) is thus
birational to \(T\) and maps onto it.  Choose any
\(z\in T_{\widehat Y}\) above \(x\).
There is an integral curve \(\Gamma\subset\widehat Y\)
through \(z\) dominating \(C_0\).  To see this, take an
irreducible component \(A\subset\pi^{-1}(C_0)\) through \(z\).
By the choice of \(x\), it dominates \(C_0\).  If
\(\dim A=d\), its fiber \(A_x\) is a proper closed subset
of dimension at most \(d-1\).  General very ample cuts through
\(z\), in number \(d-1\), cut \(A_x\) to dimension at most
zero and leave a curve component through \(z\) on \(A\).
That component cannot lie in \(A_x\), so it dominates \(C_0\).
For \(d=1\), take \(\Gamma=A\).

Let \(\iota:N\to\widehat Y\) denote the normalization map of
\(\Gamma\), followed by inclusion.
Resolve \(T_{\widehat Y}\) by a smooth projective surface
\(T_0\to T_{\widehat Y}\), choose a point \(w\in T_0\)
above \(z\), and blow it up.  Write \(T_1\) for the resulting
surface, \(e\subset T_1\) for the exceptional curve of this
blowup, and \(f:T_1\to T\) for the composite map:
\[
\begin{tikzcd}[column sep=large,row sep=large]
 T_1 \arrow[r,"\mathrm{Bl}_w"] \arrow[d,"f"'] &
 T_0 \arrow[r] &
 \widehat Y \arrow[d,"\pi"] &
 N \arrow[l,"\iota"'] \arrow[d] \\
 T \arrow[rr,hook] & & Y & C_0 \arrow[l,hook]
\end{tikzcd}
\]
The images of \(T_0\) and \(N\) in \(\widehat Y\) meet at
\(z\).  The exceptional curve \(e\), all these varieties, and
all these maps are fixed throughout the remaining argument.

\smallskip\noindent
\emph{The contradiction on the fixed surface.}
The projection formula and \(M_j\to L\) give, for one fixed
\(\delta>0\),
\begin{equation}\label{num:negative-degree}
\deg_N\iota^*\pi^*M_j\leq-\delta
\qquad\text{for all sufficiently large }j.
\end{equation}
In fact this degree equals the positive degree of
\(N\to C_0\) times \(M_j\cdot C_0\).
For each such \(j\), choose any positive multiple \(k_j\) of
\(d_j\).  Lemma~\ref{num:jets}, applied on the fixed
\(\widehat Y\) and \(\Gamma\), shows that every nonzero
pullback section
\(s'\in H^0(\widehat Y,k_j\pi^*M_j)\) of a section on \(Y\)
satisfies
\begin{equation}\label{num:uniform-point-order}
\operatorname{ord}_z(s')\geq ck_j,
\qquad c=\delta/g>0.
\end{equation}
The constant \(c\) depends only on the fixed normalized curve and
the negative-degree margin in \eqref{num:negative-degree}.

The sections on \(Y\) induce a nonzero linear system of
\(k_jf^*(M_j|_T)\), generated outside
\(f^{-1}(T\cap Z_j)\).  Nonzeroness follows from generation
at points of \(T\setminus Z_j\).

Every nonzero section in this induced system has order at least
\(ck_j\) at \(w\) before the last blowup.  Indeed, the local
ring map
\(\mathcal O_{\widehat Y,z}\to\mathcal O_{T_0,w}\)
sends \(\mathfrak m_z^b\) into \(\mathfrak m_w^b\);
apply \eqref{num:uniform-point-order} to a section inducing it.
Sections whose restriction is identically zero do not contribute
to the induced system.  Thus its fixed divisor \(F_j\) satisfies
\(\operatorname{coeff}_eF_j\geq ck_j\).
All the hypotheses of Lemma~\ref{num:fixed-square} hold with
\(D_j=M_j|_T\).  Lemma~\ref{num:fixed-square} gives a fixed \(\epsilon>0\) such
that
\[
M_j^2\cdot[T]\geq\epsilon
\qquad\text{for all sufficiently large }j,
\]
contrary to \eqref{num:special-surface-limit}.  Therefore no
curve \(C_0\) of negative \(L\)-degree exists, and \(L\) is nef.
\end{proof}
